% !TeX root = ../../../main.tex
\subsection{Broader Societal Impact}
\label{ch07:apd:broader_impact}

This work advances equitable, scalable, and transparent medical AI by introducing a framework for forecasting clinical risk in sepsis from narrative documentation.
Sepsis remains a leading cause of preventable mortality worldwide, with outcomes highly dependent on timely recognition and intervention.
In many under‑resourced healthcare settings, structured EHR data and real‑time specialist input are unavailable, leaving unstructured narratives as the primary record of patient care.
Our approach transforms such narratives into temporally structured representations, enabling systematic evaluation of forecasting methods for short‑ and long‑horizon event prediction, event ordering, and survival analysis.
By releasing the dataset, extraction pipeline, and evaluation framework, we lower barriers for developing and benchmarking temporal forecasting models in sepsis, and promote reproducibility in LLM‑based clinical information extraction.

However, societal considerations are critical.
LLM‑assisted extraction can introduce inaccuracies or hallucinations, and unvalidated forecasts risk amplifying bias or error.
Although timelines are derived from non‑identifiable public case reports, their interpretation could still reinforce historical biases if not carefully contextualized.
Forecasting results from case reports may also differ from those using routine clinical notes, limiting generalizability.
Moreover, applying these forecasting models in high‑stakes clinical decision‑making without expert oversight poses safety risks.
Responsible use requires transparent communication of dataset and model limitations, explicit uncertainty reporting, and safeguards to prevent inappropriate deployment in patient‑facing settings.

%\end{document}

